% Farbod Siahkali | One-page research internship resume
% Single-column reading order, plain text contacts, Unicode text mapping.
\documentclass[letterpaper,11pt]{article}
\usepackage[margin=0.50in]{geometry}
\usepackage{iftex}
\ifPDFTeX
  \usepackage[T1]{fontenc}
  \usepackage{lmodern}
\else
  \usepackage{fontspec}
  % Avoid ligatures and kerning gaps that can split keywords in PDF parsers.
  \setmainfont{texgyretermes-regular.otf}[
    BoldFont=texgyretermes-bold.otf,
    ItalicFont=texgyretermes-italic.otf,
    BoldItalicFont=texgyretermes-bolditalic.otf,
    Ligatures=NoCommon,RawFeature=-kern]
\fi
\usepackage{titlesec}
\usepackage{xcolor}
\usepackage{enumitem}
\usepackage[hidelinks]{hyperref}
% pdfTeX needs explicit glyph mapping; XeTeX supplies Unicode mapping.
\ifdefined\pdfgentounicode
  \input{glyphtounicode}
  \pdfgentounicode=1
\fi
\definecolor{Accent}{RGB}{0,55,105}
\hypersetup{pdftitle={Farbod Siahkali - Research Resume},pdfauthor={Farbod Siahkali}}
\pagestyle{empty}
\setlength{\parindent}{0pt}
\setlength{\parskip}{0pt}
\raggedright
\titleformat{\section}{\large\bfseries\color{Accent}}{}{0pt}{}[\titlerule]
\titlespacing*{\section}{0pt}{6pt}{4pt}
\setlist[itemize]{leftmargin=13pt,topsep=2pt,itemsep=1pt,parsep=0pt,partopsep=0pt}
\newcommand{\degree}[3]{\textbf{#1}, #2\hfill #3\par}
\newcommand{\role}[3]{\textbf{#1} | #2\hfill #3\par}
\newcommand{\paper}[3]{\item \textbf{\href{#3}{#1}}. #2.}
\newcommand{\manuscript}[2]{\item \textbf{#1}. #2.}

\begin{document}
\fontsize{10.5}{12.3}\selectfont
\begin{center}
{\fontsize{23}{25}\selectfont\bfseries Farbod Siahkali}\\[3pt]
West Lafayette, IN | +1-765-532-7185 | \href{mailto:siahkali@purdue.edu}{siahkali@purdue.edu}\\[2pt]
\href{https://scholar.google.com/citations?user=-99KSTsAAAAJ\&hl=en}{Google Scholar} | \href{https://farbod-siahkali.github.io/}{farbod-siahkali.github.io} | \href{https://www.linkedin.com/in/farbod-siahkali/}{LinkedIn: farbod-siahkali}
\end{center}
\vspace{-6pt}
Ph.D. researcher in machine learning and control developing reliable perception and planning for autonomous systems. Research spans conformal prediction, multi-sensor fusion, multi-agent trajectory optimization, and safe reinforcement learning.

\section{Education}
\degree{Purdue University}{Ph.D. in Electrical Engineering}{Expected May 2029}
\degree{Purdue University}{M.S. in Electrical Engineering, GPA: 3.83}{Expected December 2026}
\degree{University of Tehran}{B.S. in Electrical Engineering (Control), GPA: 3.92}{July 2023}

\section{Publications \& Manuscripts}
\begin{itemize}
\paper{Active Sensing and Deferred-Decision Trajectory Optimization for Robust Target Identification}
{IEEE Control Systems Letters (L-CSS), 2026. CDC 2026 presentation accepted. First author}
{https://doi.org/10.1109/LCSYS.2026.3703893}
\paper{Conformal Prediction-Based Multi-Sensor Classification}
{Neurocomputing, vol. 678, article 133182, 2026. First author}
{https://doi.org/10.1016/j.neucom.2026.133182}
\manuscript{Multi-Agent Deferred-Decision Trajectory Optimization with Anytime-Valid Target Retention for Moving Candidates}
{IEEE Robotics and Automation Letters (RA-L), under review. First author}
\paper{Coverage Guarantees for Pseudo-Calibrated Conformal Prediction under Distribution Shift}
{IEEE Signal Processing Letters, under review, 2026. arXiv:2602.14913. First author}
{https://arxiv.org/abs/2602.14913}
\manuscript{Policy Gradient with Online Parameter Estimation and Evolving CBF Action Sets}
{In preparation for submission to IEEE Control Systems Letters (L-CSS). First author}
\paper{Deep Reinforcement Learning for Epidemic Control of a Networked SIS Model}
{IEEE Conference on Control Technology and Applications (CCTA), 2025}
{https://ieeexplore.ieee.org/document/11151322}
\paper{Towards Opinion Shaping: A Deep Reinforcement Learning Approach in Bot-User Interactions}
{arXiv:2409.11426, preprint, 2024. First author}
{https://arxiv.org/abs/2409.11426}
\paper{Image-Based and Partially Categorical Annotating Approach for Pedestrian Attribute Recognition}
{SSRN, preprint, 2023}
{https://papers.ssrn.com/sol3/papers.cfm?abstract_id=4373086}
\paper{SIVD: Dataset of Iranian Vehicles for Real-Time Multi-Camera Video Tracking and Recognition}
{8th Iranian Conference on Signal Processing and Intelligent Systems (ICSPIS), 2022. First author}
{https://ieeexplore.ieee.org/document/10043932}
\end{itemize}

\section{Research Experience}
\role{Graduate Research Assistant}{Purdue University, Prof. Vijay Gupta}{Aug 2024 -- Present}
\begin{itemize}
\item Reduced prediction-set size by \textbf{10--25\%} at fixed coverage using conformal late fusion for dependent sensor outputs on synthetic, CIFAR-10, CIFAR-100, and MELD benchmarks.
\item Developed slack-inflated and source-tuned pseudo-calibration algorithms for domain adaptation on MNIST and CIFAR benchmarks. Derived coverage bounds using Wasserstein distance under label-conditional covariate shift.
\item Developed single- and multi-agent DDTO for target identification and interception using mixed-integer conic optimization, conformal prediction, sequential candidate rejection, and receding-horizon planning.
\item Developed barrier-continuation policy gradient (BCPG) for safe reinforcement learning with online parameter estimation and control barrier functions (CBFs). Evaluated adaptive cruise control and planar navigation.
\end{itemize}
\role{Research Assistant}{University of Tehran, Intelligent Networks Lab}{Oct 2023 -- Jul 2024}
\begin{itemize}
\item Implemented DDPG policies for networked epidemic control and budget-constrained opinion shaping.
\end{itemize}
\role{Research Assistant}{University of Tehran, Human \& Robot Interaction Lab}{May 2021 -- Oct 2022}
\begin{itemize}
\item Built PyTorch pipelines with YOLOv5 and OSNet for detection, re-identification, and multi-object tracking. Created CA-Duke labels for 32,000+ images and contributed to the 36,000+ image, 29-class SIVD dataset.
\end{itemize}

\section{Technical Skills}
\textbf{Languages \& Tools:} Python, MATLAB, C, Git, Linux, LaTeX.\\
\textbf{ML Libraries:} PyTorch, TensorFlow, NumPy, Pandas, SciPy, scikit-learn.\\
\textbf{Methods:} Conformal prediction, uncertainty quantification, sensor fusion, Bayesian inference, domain adaptation, mixed-integer/conic optimization, trajectory planning, multi-agent systems, safe reinforcement learning, control barrier functions, computer vision.

\section{Teaching, Service \& Honors}
Purdue ASPIRE REU research mentor (2026). University of Tehran teaching assistant (2020--2023). Reviewer: Neurocomputing, IEEE L-CSS, IEEE CDC. Best Undergraduate Project Award (2023).
\end{document}
